\section{An Exact Counterfactual What-If Interface}\label{sec:whatif}

\subsection{Model}\label{sec:model}

\paragraph{Files and row groups.} A file is identified by a fingerprint (path,
size, etag or modification time, and a BLAKE3 hash of its last 64\,KiB or,
under a stricter policy, of the whole file), not by its path, so an in-place
rewrite is detected as drift (\doc{format.md}). Each
row group receives a dense 32-bit id at the first sighting of (fingerprint,
position in file); ids are never reused, and drift tombstones the old ids
(\adr{013}). Let $\RG$ be the live row-group ids, $f(g)$ the file of row group
$g$, and $b(g)$ its compressed size from the file footer. For
$G \subseteq \RG$ let $f(G) = \{ f(g) : g \in G \}$.

\paragraph{Key specs and selections.} A key spec is $\mathsf{Exact}$ or
$\mathsf{Bloom}$ over a column list (bloom adds a target false-positive rate),
$\mathsf{Range}$ over one column, or $\mathsf{Temporal}$ over a column list
with a timestamp column. A composite key is a column list longer than one,
encoded so that byte order equals typed order and each element is
self-delimiting, so a probe on a prefix of leading elements seeks every
extension (\adr{002}, \adr{016}). For a spec $s$ and a probe $k$, the selection $S_s(k) \subseteq \RG$
is the set of row groups the index returns: for $\mathsf{Exact}$, those holding
a row whose key starts with $k$; for $\mathsf{Bloom}$, those whose filter
accepts $k$, a superset of exact membership; for $\mathsf{Range}$, those whose
minimum and maximum bound the probed value or interval; for $\mathsf{Temporal}$,
those holding a version of $k$ at its greatest timestamp (\emph{latest}) or at
the greatest timestamp not after $t$ (\emph{as-of} $t$), ties included
(\doc{format.md}). Footer statistics and embedded bloom filters alone give a
baseline selection $S_\varnothing(k)$. Let $T(k)$ be the row groups that must
be read to answer $k$ exactly. Invariant~1 requires $T(k) \subseteq S_s(k)$ on
every path. For temporal probes this needs the index-wide maximum timestamp per
key, consulted before ranking, so a superseded row is never returned even when
the newest version sits in a row group pruned for another reason
(invariant~5). Null timestamps rank last and never win (\adr{014}).

\paragraph{Runs and cost.} A recorded run $r$ carries its key set $K_r$, its
query kind (locate, latest, or as-of), and parameters such as an as-of
timestamp. Under spec $s$ it reads $G_s(r) = \bigcup_{k \in K_r} S_s(k)$, and
its cost is the vector of bytes, row groups, and files, which are the fields of
a \code{RunCost} record:
\begin{equation}\label{eq:cost}
  c_s(r) \;=\; \Big( \textstyle\sum_{g \in G_s(r)} b(g),\;\; |G_s(r)|,\;\;
  |f(G_s(r))| \Big).
\end{equation}
A physical design $D$ is a set of specs; let $D_r \subseteq D$ be the specs
whose leading columns run $r$ binds. Each selection contains $T(k)$, so their
intersection does too, and the engine intersects bitmaps over the global id
space (\adr{001}):
\begin{equation}\label{eq:design}
  G_D(r) = \textstyle\bigcap_{s \in D_r} G_s(r), \qquad
  G_D(r) = G_\varnothing(r) \text{ if } D_r = \emptyset .
\end{equation}

\subsection{Why the cost is exact}\label{sec:exact}

\emph{Design property.} $c_D(r)$ is computed without an estimated quantity.
$K_r$ is recorded. Each $S_s(k)$ is a deterministic function of the files: key
membership for exact and temporal specs, filter bits for bloom specs (fixed by
the canonical encoding and the hash function; \adr{008}, \adr{016}), and footer
bounds for range specs. $b(g)$ is read from the footer. There is no cardinality
estimate, cost model, or sampling, and the same workload, manifest, and
proposal yield a byte-identical Replay (invariant~3). The predicted quantity is
also the one the plan prescribes: a consumer that executes the scan plan reads
the row groups in $G_D(r)$ and applies the residual predicate (\adr{011}).

Three boundaries keep the claim precise, and the paper's use of ``exact'' is
narrowed to them. (1)~\emph{I/O as accounted by row-group metadata}: $b(g)$ is
a row group's total compressed size, a consumer that projects columns reads a
subset of it, and per-column chunk sizes are not yet in the statistics type
(\doc{rule-catalog.md}), so projection savings are outside the exact cost, as
are reads of footers, segments, and the manifest. (2)~\emph{Files that exist}:
a hypothetical key spec over existing files is exact, but a hypothetical layout
changes the row-group statistics, which do not exist until the layout is
written. Its cost is exact once the layout is materialized, and until then no
number is reported, because \code{RunCost} has no field for an estimate
(\doc{rule-catalog.md}, \code{scattered\us time\us column}). (3)~\emph{Scans whose
key sets are recorded or derived from the run's inputs}, such as trigger keys
and parameters: costs above the scan (CPU, memory, latency) and scans whose key
sets depend on unrecorded intermediate results are measured by an attached
executor through the planned \code{Scorer} trait, not replayed
(\doc{evolve.md} \S10).

\subsection{Computing a counterfactual}\label{sec:counterfactual}

A counterfactual asks what a recorded run would have read under a design that
differs from the current one. Range specs and the statistics baseline need only
footers. For an exact, bloom, or temporal spec that has not been built, replay
constructs a transient in-memory segment from a projected read of the key
columns (and the timestamp column for temporal specs), one segment at a time
(\doc{work-packages.md}, \adr{018}). The read is paid once per column set and
shared by every run and every candidate over those columns; each run then costs
one lookup per key, one bitmap union, and one sum.

Prefix seek models subsumption: an exact spec on $(a, b)$ answers a probe on
$a$ alone, so for runs keyed by $a$ it replaces a spec on $a$. Nulls are handled
per tuple (\adr{014}, as amended): a tuple is indexed when its leading column is
non-null, a null in a trailing column is encoded with the null marker so the
tuple stays prefix-seekable, and a probe containing a null matches nothing, as
in SQL's three-valued logic. The composite's prefix selection therefore equals
the single-column selection, including rows whose trailing columns are null, so
dropping the single-column spec preserves invariant~1.

\subsection{Index structures and complexity}\label{sec:segments}

Segments are immutable and content-addressed, and each indexes one key spec
over a range of row-group ids (\doc{format.md}). An exact segment stores sorted
key blocks, a split-block bloom filter per block~\cite{parquetbloom}, and for
each key a roaring bitmap~\cite{roaring} over global row-group ids
(\adr{001}). A temporal segment adds, per key, a chain of (timestamp, row group)
pairs and the key's maximum timestamp. A bloom segment stores one split-block
filter per row group, the shape Parquet embeds, and records whether a filter
hashes canonical bytes or Parquet's plain encoding so that a probe is encoded
to match (\adr{008}, amendment). A range segment stores a minimum and maximum
per row group. One canonical byte encoding per Arrow type serves ordering,
bloom hashing, and the row hashing of Section~\ref{sec:cert} (\adr{016}).
Adding a key spec builds segments over every row group and records completeness
in the manifest; a query on an incomplete spec fails rather than returning a
partial answer (\adr{012}).

For a segment with $n$ distinct keys, an exact lookup of key set $K$
binary-searches the block index, probes one block's filter, and searches that
block's key directory only if the filter passes, for $O(|K| \log n)$
comparisons, while holding one block index and the bitmaps of the keys asked
(\adr{018}). The union is proportional to the bitmaps' containers and the byte
sum is $O(|G_D(r)|)$. Building a transient segment is one projected pass over
the key columns, with memory bounded by one segment (\doc{format.md}).
The design budget for one replay of one run under one candidate is 5\,ms on the
medium synthetic dataset (\doc{evolve.md} \S12); the harness asserts it, and it
is a budget, not a measurement.
